\section{Related Work}
\label{sec:related_works}

\textbf{Structured action representations.}
ACT and Diffusion Policy predict action chunks to model temporally extended behavior \cite{zhao2023act,chi2023diffusion}. VLAs combine such action prediction with pretrained visual and language representations \cite{black2024pi0,shukor2025smolvla}; FAST compresses action sequences with frequency-domain tokens \cite{pertsch2025fast}. Several recent methods use splines. BEAST encodes action sequences as B-spline tokens \cite{zhou2025beast}. B-spline Policy predicts continuous curves and accelerates execution by temporal rescaling \cite{han2026bsp}. Spline Policy studies continuous outputs across policy families and their use in resampling, editing, and downstream control \cite{tian2026spline}. Our contribution concerns event-duration supervision and the interaction between a structured action head and executable language. Continuous trajectories and temporal resampling are shared with this prior work.

\textbf{Language and long-horizon execution.}
Large VLAs combine heterogeneous training data and semantic prediction to execute extended tasks \cite{black2025pi05}. Long-VLA uses phase-aware inputs \cite{fan2025longvla}, while tool-aligned VLAs separate high-level reasoning from policy invocation \cite{lei2026tools}. We instead study a single policy conditioned on successive clauses, with the physical training trajectories held fixed. Counterfactual instruction studies expose cases where visual or positional shortcuts dominate language \cite{glossop2025cast,fang2026counterfactual}. They motivate testing the requested behavior directly. Our order intervention therefore measures the first completed subgoal separately from final task success; changing the reward under an incompatible prompt alone would not establish successful steering.

\textbf{Execution timing.}
PACE selects action-chunk execution horizons from low-speed transitions at inference \cite{nie2026pace}. Inference-time policy steering uses human interaction to guide generated behavior \cite{wang2025itps}. Here, demonstration events supervise segment duration, and a decoder controls its execution time. The clause scheduler and motion clock act at different scales: one chooses the active subgoal; the other determines how a predicted segment is sampled.
